\section{Characterizing and Organizing the Multi-Fidelity Simulation Data}

\subsection{Existing Training Dataset}

\begin{table}[ht]
    \centering
    \caption{Comparative Analysis of BRUTUS Training Datasets. The new dataset provides a dramatic 8-fold to 17-fold increase in the temporal resolution of the ground-truth trajectories.}
    \label{tab:brutus_compare}
    \begin{tabular}{@{}lllllll@{}}
    \toprule
    \thead{Parameter} & \thead{Dataset/\\Number of Files} & \thead{IC Type} & \thead{tolerance} & \thead{Word\\Length} & \thead{dt\\out} & \thead{Temporal Density\\Increase} \\ \midrule
    \textbf{Original}  & Set 1a (4857)           & Hierarchical     & $1\text{e-}15$     & 130                  & 0.08             & Baseline (1x)                      \\
                       & Set 1b (4857)          & Standard         & $1\text{e-}13$     & 118                  & 0.17             & Baseline (1x)                      \\
    \textbf{New}       & Set 2a (2749)          & Hierarchical     & $1\text{e-}13$     & 130                  & \textbf{0.01}    & \textbf{8x}                        \\
                       & Set 2b (2565)          & Standard         & $1\text{e-}13$     & 130                  & \textbf{0.01}    & \textbf{17x}                       \\ \bottomrule
    \end{tabular}
\end{table}

For \textbf{Set 1a Hierarchical}, 5000 raw ICs were generated; \textbf{Set 1b Standard}, 4000. For both types of ICs, \textbf{Set 2a Hierarchical} and \textbf{Set 2b Standard}, the same number of ICs were used.
However, these are just the raw ICs. The actual number of IC files left are subject to the prescreening result.
Furthermore, the number of BRUTUS output files are even less due to the computational limitation. The final number of BRUTUS output files are in the above table – for each set,
it is the same number for both the IC files and the BRUTUS output files. In total there were 15,028 pairs of IC and BRUTUS output files generated for training data,
ready to be parsed in the next stage for preparing this training data.


\subsection{Multi-Fidelity Dataset and Physical Invariants}

A robust training dataset is crucial for PINNs on the three-body problem.  We utilize two complementary sources of simulated trajectories: (1) High-fidelity "hierarchical" simulations, and (2) Lower-fidelity "quasi-random" simulations.  The hierarchical dataset consists of stable three-body configurations (a tightly bound binary with a distant tertiary companion) generated with a small time step ($dt \approx 0.08$) for high precision.  The quasi-random dataset covers a broad variety of chaotic configurations generated by random initial conditions (using Sobol quasi-random sampling for good coverage) and integrated with a larger time step ($dt \approx 0.17$) for efficiency.  All simulations are done with the Brutus N-body integrator - a high-precision solver - but the second dataset, being more random, is considered "low-fidelity" due to slightly lower precision and more chaotic behavior.  Each simulation's initial conditions determine the system's characteristics. We compute five key physical invariants or properties from the initial state of each simulation, which serve as quantitative difficulty metrics: 

\begin{itemize}
    \item \textbf{Total Mechanical Energy ($E$):} The sum of kinetic and potential energy of the system.  $E<0$ indicates a bound (stable) system, 
    with more negative $E$ meaning a deeper gravitational well (tighter binding).  Values near zero indicate marginally bound or unbound 
    systems which tend to be more chaotic.  (Bound systems satisfy $E<0$; as a guideline, "Stage 0" easiest cases have $E<-0.5$ in normalized units.)


    \item \textbf{Total Angular Momentum Magnitude ($L$) and $L/L_{\max}$:} $L$ is the magnitude of the total angular momentum vector. 
    We often normalize it by $L_{\max}$, the theoretical maximum angular momentum for the system's energy (which corresponds to 
    all mass moving on a circular orbit for that energy). A high $L/L_{\max}$ (close to 1) means near-circular orbits (stable); 
    a low ratio means elongated or radial orbits that lead to closer encounters and chaos. (For example, Stage 0 requires $L/L_{\max}>0.5$ - 
    reasonably high angular momentum.)



    \item \textbf{Virial Ratio (Q):} The ratio $2T/|V|$ of twice the total kinetic energy to the absolute potential energy at $t=0$.  
    For bound stable systems in virial equilibrium, $Q \approx 1$. Significant deviations (too low or high) imply an imbalance (e.g. 
    too much kinetic energy or too much potential energy) that can lead to expansion or collapse.  Moderate ranges around 1 indicate 
    stability (Stage 0 uses $0.7 \le Q \le 1.3$ as a criterion). 



    \item \textbf{Minimum Initial Separation ($s_{\min}$):} The smallest distance between any two bodies in the initial configuration. 
    A larger $s_{\min}$ means the bodies are well separated initially, reducing the chance of immediate close encounters.  
    Very small initial separations (two bodies almost colliding at $t=0$) indicate extremely challenging cases that can cause 
    numerical difficulties and chaotic scattering.  Stage 0 imposes $s_{\min}>2$ (in code units) to avoid any incipient close approaches. 
\end{itemize}

All data is normalized to dimensionless code units for training (e.g., positions and distances scaled by a characteristic length, velocities by a characteristic orbital velocity) to keep values typically in $[-1,1]$.  A consistent normalization is applied across the entire dataset, which aids network training by ensuring inputs and outputs are on comparable scales. 










\subsection{Multi-Stage Data Stratification by Difficulty (Revised to handle the issue that virial ratio is not preserved for quasi-random IC sampling)}

\authcomment{I have a suggestion: how about we give the entire IC files a scan using some algorithms and analyze the gathered data from there,
to first get an insight on what the data looks like, and then we can
decide on the criteria for each stage? This will be a good starting point considering that simply designing the criteria based on theoretical values
might not be the best way to go espesically since my IC sampling method is complicated – virial ratio is not preserved for quasi-random IC sampling.}

With these invariants computed for all simulated trajectories, we can quantitatively stratify the data from "easy" (
stable, regular motion) to "hard" (chaotic, complex outcomes).  The simulations are binned into a multi-stage curriculum 
based on thresholds of the invariants.  The binning strategy (summarized from Table 1 of the original document) is as follows:


\begin{itemize}

    \item \textbf{Stage 0: Easiest, highly stable systems.}  Data source: high-fidelity hierarchical simulations only (Set 1a). 
    \begin{itemize}
        \item \textbf{Criteria:}
        \begin{itemize}
            \item Total Mechanical Energy: $E < -0.5$
            \item Angular Momentum Ratio: $\frac{L}{L_{\max}} > 0.5$
            \item Virial Ratio: $0.7 \leq Q \leq 1.3$
            \item Minimum Initial Separation: $s_{\min} > 2$
        \end{itemize}
        \item \textbf{Description:} Tightly bound binaries with near-virial equilibrium and well-separated initial positions. Ideal for initial training due to smooth trajectories.
    \end{itemize}
    
    
    
    \item \textbf{Stage 1: Still stable but slightly more complex.} Data source: remaining hierarchical simulations, 
    divided into two sub-stages to introduce one complexity at a time (Set 1a).  

    \begin{itemize}
        \item \textbf{Stage 1a: Moderately Stable (Energy Relaxed)}
        \begin{itemize}
            \item \textbf{Criteria:}
            \begin{itemize}
                \item $-0.5 \leq E < -0.1$
                \item $\frac{L}{L_{\max}} > 0.5$
                \item $0.7 \leq Q \leq 1.3$
                \item $s_{\min} > 2$
            \end{itemize}
            \item \textbf{Description:} Systems with moderately reduced binding energy while maintaining orbital stability.
        \end{itemize}

        \item \textbf{Stage 1b: Moderately Stable (Angular Momentum Relaxed)}
        \begin{itemize}
            \item \textbf{Criteria:}
            \begin{itemize}
                \item $E < -0.5$
                \item $0.1 \leq \frac{L}{L_{\max}} < 0.5$
                \item $0.7 \leq Q \leq 1.3$
                \item $s_{\min} > 2$
            \end{itemize}
            \item \textbf{Description:} Systems with moderately reduced angular momentum (eccentric orbits) 
            but strong binding energy.
        \end{itemize}
    \end{itemize}
    
    
        
    \item \textbf{Stage 2 (High-Resolution Refinement of Stable Orbits):} 
    Data source: High-resolution hierarchical simulations (Set 2a).
    \begin{itemize}
        \item \textbf{Criteria:} Union of Stage 0 and Stage 1 criteria
        
        \item \textbf{Description:} 
        This is the first critical revision. 
        After mastering stable orbits in a low-resolution format, 
        the model is now tasked with refining its predictions using the high-resolution BRUTUS output of the \emph{same} physical 
        systems (with the same hierarhical IC files used for \textbf{Stage 0} and \textbf{Stage 1}). 
        Because the underlying physics is simple and already learned, 
        the network can focus its capacity entirely on fitting the more detailed temporal signal without the compounding 
        difficulty of chaotic dynamics. 
    \end{itemize}


    
    \item \textbf{Stage 3: Moderately Chaotic (Low-Fidelity)}  Data source: low-fidelity quasi-random simulations (Set 1b).
    \begin{itemize}
        \item \textbf{Criteria:}
        \begin{itemize}
            \item $E \leq E^{\mathrm{LF}}_{25\%}$ \hspace{0.5em} (top 25\% most bound in LF dataset)
            \item $\frac{L}{L_{\max}} \geq L^{\mathrm{LF}}_{75\%}$ \hspace{0.5em} (top 25\% highest angular momentum)
            \item $0.9 \leq Q \leq 1.1$
        \end{itemize}
        \item \textbf{Description:} Least chaotic systems from low-fidelity data. Introduces early chaos while maintaining relative stability.
    \end{itemize}
    
    
    \item \textbf{Stage 4: Fully chaotic regimes, introduced in steps.} 
    Data source: Remaining low-fidelity data (Set 1b).  

    
    \begin{itemize}
        \item \textbf{Stage 4a: Chaotic (Long Predictability Horizon)}
        \begin{itemize}
            \item \textbf{Data Source:} Remaining low-fidelity simulations
            \item \textbf{Criteria:}
            \begin{itemize}
                \item $\tau_L > 5$ \hspace{0.5em} (Lyapunov time $>$ 5 code units)
            \end{itemize}
            \item \textbf{Description:} Systems with moderate chaos and relatively long predictability horizons.
        \end{itemize}
        
        \item \textbf{Stage 4b: Highly Chaotic (Short Predictability Horizon)}
        \begin{itemize}
            \item \textbf{Data Source:} Remaining low-fidelity simulations
            \item \textbf{Criteria:}
            \begin{itemize}
                \item $2 < \tau_L \leq 5$
            \end{itemize}
            \item \textbf{Description:} Systems with strong chaos and limited predictability.
        \end{itemize}
    \end{itemize}


    \item \textbf{Stage 5: High-Resolution Chaotic Orbits}
    Data source: High-resolution chaotic simulations (Set 2b).
    
    \begin{itemize}
        \item \textbf{Criteria:} Same as Stage 4
        \item \textbf{Description:} High-temporal-resolution trajectories of chaotic systems. Focuses on resolving close encounters.
    \end{itemize}


    \item \textbf{Stage 6: Near-collision and extreme cases.} 
    Data source: All remaining data; HF+LF 
    
    \begin{itemize}  
        \item \textbf{Criteria:}
        
        $s_{\min} < 0.5$ \textbf{OR} $|E| < 0.01$
  
        \item \textbf{Description:} Near-collision scenarios and borderline unbound systems.
    \end{itemize}

\end{itemize}



Crucially, each subsequent stage's training set includes all data from previous stages (curriculum replay) plus the new subset.  
This means by the time we reach Stage 5 or 6, the training data is the entire dataset.  
However, because the model has been gradually brought along this path, it is not learning from scratch on this complex set;  
it has a strong foundation from simpler cases. The staged curriculum defined above is empirically designed to ensure a smooth 
progression from stable to chaotic dynamics, aligning with human intuition of difficulty - initial stages cover basic Keplerian 
motion and conservation laws, and later stages add perturbations, resonances, and eventually close encounters and escapes.  
This invariant-based stratification is mathematically sound in that each invariant correlates with aspects of the equations 
of motion's behavior (stability, boundedness, etc.), providing a principled way to rank problem instances by difficulty. 


\begin{comment}
\subsection{Assessing Data Suitability and Generation Methodology}

\textbf{Coverage:} The combination of hierarchical and quasi-random simulation data yields a comprehensive coverage of the scenario space.  The hierarchical dataset ensures that the model sees clean, noise-free examples of stable orbits (important for learning baseline gravitational dynamics accurately).  The quasi-random dataset contributes diversity including highly chaotic and rare events which the model must eventually handle 3 5. Importantly, because the chaotic data is integrated with a somewhat larger timestep, it contains slight numerical noise;  introducing it only after the model has learned a strong "physics prior" on the noiseless data helps prevent the network from fitting integration errors or spurious artifacts 5.

\textbf{Data Quality Control:} The data generation pipeline itself (detailed in the companion methodology document) was carefully constructed to ensure high-quality training data despite the chaotic nature of the problem.  A multi-stage filtering pipeline was employed 6 7:

\begin{itemize}
    \item In Stage 1 of generation, initial conditions are sampled via two methods: a hierarchical system generator (producing stable configurations by design) and a Sobol low-discrepancy random sampler (to uniformly cover the space of more chaotic configurations) 8 9. This yields a diverse pool of candidate initial conditions. 
    \item \textbf{Stage 2 (Pre-screening):} Each candidate is subjected to a short, low-precision integration using a fast leapfrog integrator 10. This inexpensive test identifies blatantly unstable cases before investing in a costly high-precision integration.  Specifically, any IC that shows $>1\%$ energy drift in a very short time or that has a dangerously close approach ($<0.01$ in distance) during this trial is rejected 11 12. This step weeds out configurations that would have failed or required extremely small time-steps in the precise Brutus integrator, thereby saving a huge amount of computation 10 13. The result is that only "promising" initial conditions - those likely to conserve energy well and not immediately go singular - are passed on. 
    \item \textbf{Stage 3 (Adaptive high-precision simulation):} The surviving ICs are then integrated with Brutus, but not all with the same settings.  Based on each system's properties (e.g. expected timescales, closeness of encounters), the integrator's parameters (like error tolerances or maximum step sizes) are tuned for each regime 14. For example, a system predicted to have very close passes will be simulated with more stringent error controls (smaller time-steps, higher precision) than a system that seems benign.  This adaptive parameterization ensures that we do not overspend computational effort on easy cases while still maintaining accuracy for hard cases 15. The net effect is a hierarchically generated, quality-checked dataset that is rich enough for learning but does not waste resources on untrainable or junk data. 
\end{itemize}

Thanks to this pipeline, the datasets are well-suited for curriculum learning.  The easier stages (0-1) are drawn from the meticulously stable hierarchical runs (which passed all checks and have small step sizes), providing trustworthy ground truth for the PINN to latch onto.  The chaotic data introduced later has also passed basic energy-conservation and resolution checks, meaning the PINN won't be asked to learn from inconsistent or nonsensical training targets.  The five invariants used to define the curriculum stages are indeed the same metrics used during data generation (e.g., virial ratio and minimum separation were used as filters) - this consistency reinforces that the curriculum progression is grounded in physical measures of complexity.  Overall, the data generation strategy is effective and we recommend continuing with this hybrid approach for future data expansion.  It achieves a wide coverage without requiring higher precision or significantly more simulations.  Any gaps in the difficulty spectrum can be filled by targeted new simulations in an active learning loop (see Section 4.3) rather than by overhauling the sampling strategy.  In short, the current datasets and stage definitions are sufficiently comprehensive and well-structured for the learning task, striking a balance between fidelity, diversity, and manageability. 
\end{comment}


